\chapter*{Introduction générale}
\addcontentsline{toc}{chapter}{Introduction générale}
\markboth{Introduction générale}{}

À l'ère de la transformation numérique, l'Industrie~4.0 marque une évolution majeure des secteurs industriels et des laboratoires. Elle repose sur l'intégration des technologies numériques, de l'Internet des objets (\ac{IoT}), des systèmes d'information et de l'automatisation des processus \citep{kagermann2013,lasi2014}. Cette transformation améliore la performance opérationnelle, renforce la qualité des produits et des analyses, assure une traçabilité complète des opérations et facilite la prise de décision grâce à l'exploitation des données en temps réel.

Les laboratoires d'analyses sont directement concernés. Soumis à des exigences strictes de fiabilité et de traçabilité, notamment celles de la norme \ac{ISO/CEI 17025} \citep{iso17025}, ils produisent chaque jour un grand volume de données : demandes, échantillons, mesures, résultats et rapports. Lorsque ces données sont saisies à la main et dispersées sur plusieurs supports, les risques d'erreur augmentent, les délais s'allongent et la reconstitution de l'historique devient difficile.

C'est dans ce contexte que s'inscrit ce projet de fin d'études, réalisé au sein de BQube ITs, entreprise tunisienne spécialisée dans les solutions numériques pour l'industrie et les laboratoires, en particulier les systèmes \ac{MES} et \ac{LIMS}. Le projet comporte deux volets complémentaires qui couvrent les deux niveaux de la transformation numérique d'un laboratoire :
\begin{itemize}
  \item \textbf{le niveau du système d'information}, au Centre Technique de l'Agro-Alimentaire (\ac{CTAA}) : participation au déploiement d'un \ac{LIMS} (paramétrage, étiquetage par codes-barres et codes \ac{QR}, formation des utilisateurs), puis conception et développement d'un module de traçabilité ;
  \item \textbf{le niveau des équipements}, aux Laboratoires Pharmaghreb : étude de la connectivité de dix équipements de laboratoire, puis conception, réalisation et test d'un prototype d'acquisition qui récupère automatiquement les mesures d'une balance par liaison \ac{RS-232} et les transmet par \ac{Wi-Fi}. Ce prototype est fonctionnel et attend son déploiement.
\end{itemize}

La problématique générale peut se formuler ainsi : \emph{comment digitaliser les processus analytiques d'un laboratoire, de la gestion des échantillons jusqu'aux machines connectées, tout en garantissant la fiabilité et la traçabilité des données ?}

Ce rapport s'organise en quatre chapitres :
\begin{itemize}
  \item le \textbf{premier chapitre} présente le cadre général du projet : l'entreprise d'accueil, les deux cas d'étude, l'analyse de l'existant, la problématique, la méthodologie et la planification ;
  \item le \textbf{deuxième chapitre} est consacré à l'analyse et à la spécification des besoins ;
  \item le \textbf{troisième chapitre} détaille la conception du module de traçabilité et du prototype de machine connectée ;
  \item le \textbf{quatrième chapitre} décrit la réalisation : environnement de travail, déploiement du \ac{LIMS}, interfaces développées, tests, formation des utilisateurs et réalisation du prototype.
\end{itemize}
Une conclusion générale dresse le bilan du travail et ouvre des perspectives. L'annexe~\ref{annexe:gains} détaille la méthode d'estimation des gains économiques.
